\section{The original residue pairing on the product base and its
diagonal
pullback}\label{the-original-residue-pairing-on-the-product-base-and-its-diagonal-pullback}

24 September 2026. Independent derivation RPD0--RPD10.

\subsection{RPD0. The exact continuation and its
sources}\label{rpd0.-the-exact-continuation-and-its-sources}

\href{https://github.com/KokunoYumeto/zeta-function-research-reader/blob/main/workbenches/splitzero-tandem/continuations/20260925-full-source-tensor-and-original-divisor/counterfactual/FINITE_TOPOLOGY_DUALIZING_COMPLEX.md}{FTD0--FTD10}
constructed the right adjoint on the actual Connes--Consani three-point
topology, with all stalks, restrictions, and the continuous/algebraic
dual comparison. It proved that the existing nonzero global map into the
character-twisted dual shifted by \(-2\) cannot be the global sections
of a degree-zero same-site sheaf map with that target. Following the
amended forward rule and the whole established corpus, the next
operation here is the product-site pairing, followed by an explicit
diagonal pullback. Neither operation changes the original arithmetic
quotient.

The current corpus rules (private construction record; not distributed),
\href{https://github.com/KokunoYumeto/zeta-function-research-reader/blob/36a82ecca98addb8a98b48204e349737856c7223/workbenches/splitzero-tandem/continuations/20260924-cohomology-weight-and-character-lifts/original-character-lifting/tau_lifting_weights_20260924/SOURCE_CC_DOUBLE_PULLBACK.md}{DCP0--DCP12},
\href{https://github.com/KokunoYumeto/zeta-function-research-reader/blob/36a82ecca98addb8a98b48204e349737856c7223/workbenches/splitzero-tandem/continuations/20260924-cohomology-weight-and-character-lifts/original-character-lifting/tau_lifting_weights_20260924/SOURCE_COEFFICIENT_GLUE.md}{CGS0--CGS10},
and the full user arguments recorded by FTD0 remain controlling. Source
\(Z_0,Z_1/\tau,Z_2\) keep their supplied meanings. No source addition,
numerical weight or parity is assigned to primitive \(\tau\).

The human geometric source is Alain Connes and Caterina Consani,
\href{https://arxiv.org/abs/0903.2024v3}{\emph{Schemes over
\(\mathbb F_1\) and zeta functions}, arXiv:0903.2024v3, §5}, with
original author-source coverage recorded in DCP0. The analytic source
framework is Ralf Meyer,
\href{https://arxiv.org/abs/math/0412277v3}{arXiv:math/0412277v3}, with
the full original-zeta comparison proved in
\href{https://github.com/KokunoYumeto/zeta-function-research-reader/blob/36a82ecca98addb8a98b48204e349737856c7223/workbenches/splitzero-tandem/continuations/20260924-cohomology-weight-and-character-lifts/original-character-lifting/tau_lifting_weights_20260924/ORIGINAL_MELLIN_SPECTRAL_SYNTHESIS.md}{OMS}.
Deligne's geometric weight argument remains
\href{https://www.numdam.org/item/PMIHES_1980__52__137_0/}{\emph{La
conjecture de Weil. II}, §§3.6.1--3.6.3}, whose actual cross was read
and reconstructed in ORE7. No geometric purity statement is transferred
by assumption.

The actual residue pairing, including its original function, continuity,
full multiplicity, descent and real action, is the existing
\href{https://github.com/KokunoYumeto/zeta-function-research-reader/blob/36a82ecca98addb8a98b48204e349737856c7223/workbenches/splitzero-tandem/continuations/20260924-cohomology-weight-and-character-lifts/original-character-lifting/tau_lifting_weights_20260924/GLOBAL_ORIGINAL_ZETA_RESIDUE_DUALITY_INDEPENDENT.md}{GZR0--GZR9}.
This proof constructs its new sheaf-level product realization and all
requested comparison maps, rather than claiming the original pairing is
newly derived.

\subsection{RPD1. Original spaces and full
contour}\label{rpd1.-original-spaces-and-full-contour}

Use the actual finite base \[
Y=\{c_+,c_-,\eta\},\qquad
U_\pm=\{c_\pm,\eta\},\quad U_\eta=\{\eta\}.
\] The faithful original coefficient sheaf is \[
F=\Omega_{\rm full}
=(W_+\xrightarrow{r_+}A\xleftarrow{r_-}W_-),
\quad
W_\pm=(S\oplus\mathbb C^2)\oplus V_\pm^{\rm extra}.
\tag{RPD1.1}
\] Here \[
S=\{h\in\mathcal S(\mathbb R;\mathbb C):
h(-v)=h(v),\ h(0)=0,\ \int_{\mathbb R}h=0\},
\] \[
A=\{b\in C^\infty(\mathbb R_{>0}):
\sup_{u>0}(u^N+u^{-N})|(u\partial_u)^jb(u)|<\infty
\quad(N,j\ge0)\}.
\] The restrictions are \[
r_+(h,c,w)=\Sigma h,\quad r_-(h,d,w)=R\Sigma h,\quad
\Sigma h(u)=2\sum_{n\ge1}h(nu),\quad Rb(u)=u^{-1}b(u^{-1}).
\tag{RPD1.2}
\] Every endpoint and every extra coordinate has zero restriction and
remains in its chart space. The closed image theorem gives \[
r_+(W_+)=r_-(W_-)=J=\Sigma S,\qquad Q=A/J,\qquad \pi:A\to Q.
\] The original Čech complex is \[
D=[P:=W_+\oplus W_-\xrightarrow{d=r_+-r_-}A]
\quad(\deg=0,1).
\tag{RPD1.3}
\]

Write \(\mathcal M_0a(s)=\int_0^\infty a(u)u^s\,du/u\). Define the
original raw pairing by \[
B_A(a,b)=\frac1{2\pi i}
\left(\int_{2-i\infty}^{2+i\infty}
-\int_{-1-i\infty}^{-1+i\infty}\right)
\frac{\mathcal M_0a(s)\mathcal M_0b(1-s)}{\zeta(s)}\,ds.
\tag{RPD1.4}
\] Both edges are upward. This is the original denominator \(\zeta(s)\).
The original comparison with an earlier centered representative is
\(\mathcal T k(u)=2u^{-1/2}k(u)\) and
\(\mathcal M_0\mathcal T k=2\mathcal M k\); consequently the
corresponding two-input formula on \(\mathcal T k,\mathcal T\ell\)
retains the factor4. No centered replacement is used here.

GZR proves, with
\(b_{2,1}(F)=\sup_{|\sigma|\le2,t\in\mathbb R}(1+|t|)|F(\sigma+it)|\),
\[
|B_A(a,b)|\le
\frac{\zeta(2)(1+4\pi^2)}{\pi}\,
b_{2,1}(\mathcal M_0a)b_{2,1}(\mathcal M_0b).
\tag{RPD1.5}
\] The integrals converge separately. Its original-source division proof
gives \[
B_A(J,A)=B_A(A,J)=0,\qquad
B_A(a,b)=B_\zeta(\pi a,\pi b).
\tag{RPD1.6}
\] The endpoint values are \(\zeta(0)=-1/2\) and a simple zero of
\(1/\zeta\) at1. The trivial zeros are outside the two stated edges;
their original inverse-Mellin residues are \[
\frac{F(-2r)}{\zeta'(-2r)},\qquad
\zeta'(-2r)=(-1)^r2^{-2r-1}\pi^{-2r}(2r)!\zeta(1+2r).
\tag{RPD1.7}
\] They remain in the full return theorem invoked for RPD1.6. The
functional equation uses the complete multiplier \[
\zeta(s)=\chi_\zeta(s)\zeta(1-s),\qquad
\chi_\zeta(s)=\pi^{s-1/2}
\frac{\Gamma((1-s)/2)}{\Gamma(s/2)}.
\tag{RPD1.8}
\] These factors are not replaced by a completed denominator.

The diagonal real action is \(T_ab(u)=b(u/a)\) with the original
compatible chart actions. Direct substitution gives \[
B_A(T_aa,T_ab)=a\,B_A(a,b).
\tag{RPD1.9}
\] The two factors acquired by the Mellin integrand are \(a^s\) and
\(a^{1-s}\), whose product is the actual character
\(\chi_{\rm dil}(a)=a\).

\subsection{RPD2. Construct the product sheaf morphism on all nine
points}\label{rpd2.-construct-the-product-sheaf-morphism-on-all-nine-points}

The product \(Y\times Y\) has nine points and minimal open
\(U_x\times U_y\) at \((x,y)\). Define the external tensor sheaf \[
(F\boxtimes F)_{(x,y)}=F_x\otimes_kF_y,\qquad k=\mathbb C,
\tag{RPD2.1}
\] with every restriction the tensor product of the original
restrictions. This is the algebraic tensor over the stated receiving
field. It is the actual external tensor of sheaves: the inverse-image
stalks from both projections are \(F_x,F_y\), and sheaf tensor is
computed stalkwise.

Let \(g=(\eta,\eta)\) and \(j_g:\{g\}\hookrightarrow Y^2\). The
extension by zero \[
K=j_{g!}(\chi_{\rm dil}k)
\tag{RPD2.2}
\] has stalk \(\chi_{\rm dil}k\) at \(g\), zero at every other point,
and zero maps on every arrow with zero source. This follows directly
from the diagram extension-by-zero adjunction: a map from this diagram
is determined by its generic component.

Define \(\beta:F\boxtimes F\to K\) to be \(B_A\) at \(g\) and zero at
the other eight points. It is a sheaf morphism. To check every
naturality square ending at \(g\), any incoming restriction from a point
other than \(g\) has image in \(J\otimes A\), \(A\otimes J\), or
\(J\otimes J\), according to whether its first, second, or both
coordinates are closed. Equation RPD1.6 makes \(\beta_g\) zero on each
such image. All other squares have zero target components. Thus \[
\boxed{\beta:F\boxtimes F\longrightarrow j_{(\eta,\eta)!}\chi_{\rm dil}k}
\tag{RPD2.3}
\] is constructed on the entire product topology.

It is equivariant for simultaneous dilation on both factors by RPD1.9,
and therefore for every original prime action. No one-dimensional
character for two independently chosen dilations is asserted: the proved
character covariance is for the stated diagonal group action.

The topology is also explicit. Each original stalk is Fréchet; put the
projective locally convex tensor topology on each algebraic tensor in
RPD2.1. The original restrictions are continuous, and RPD1.5 makes the
generic functional continuous. It extends uniquely to the completed
projective tensor because \(k\) is complete. Thus the same component
formulas define a continuous morphism on that completed diagram as well.
All cohomology tensor identities below refer to the algebraic sheaf
calculation unless stated otherwise; no completed-tensor exactness
theorem is silently assumed.

\subsection{RPD3. Ordered product Čech complex and trace in
degree2}\label{rpd3.-ordered-product-ux10dech-complex-and-trace-in-degree2}

Use the plus/minus ordering in each factor. The original length-one
projective resolution of \(k_Y\) is FTD2.1, with the map at the generic
point \(t\mapsto(-t,t)\). Its external tensor with itself is a
length-two projective resolution of \(k_{Y^2}\). Exactness can be
checked stalkwise: each factor is a split exact resolution of a
one-dimensional vector space, and the tensor total complex has the same
degree-zero cohomology and zero higher cohomology. Its projective terms
represent the nine stalk evaluations, so it computes derived global
sections for every sheaf on \(Y^2\).

Applying this resolution to \(F\boxtimes F\), with the ordinary Koszul
identification of the two Hom complexes, gives the ordered product
complex \[
C^0=P\otimes P,\quad
C^1=(A\otimes P)\oplus(P\otimes A),\quad
C^2=A\otimes A,
\tag{RPD3.1}
\] \[
d_C^0(x\otimes y)=(dx\otimes y,\ x\otimes dy),
\] \[
d_C^1(a\otimes y,\ x\otimes b)
=-a\otimes dy+dx\otimes b.
\tag{RPD3.2}
\] Thus \(C=D\otimes D\) with differential
\(d\otimes1+(-1)^{\deg}1\otimes d\). The two terms of \(d_C^1d_C^0\) are
\(-dx\otimes dy\) and \(+dx\otimes dy\), so their sum is zero.

The degree-two orientation is part of this formula. In the literal Hom
of the external projective resolution, the identification with
Hom-tensor has sign \((-1)^{ij}\) for degrees \(i,j\). At \(i=j=1\) it
is \(-1\). This is the Koszul sign that gives the degree-one
differential in RPD3.2. After this specified identification the top
coordinate \(A\otimes A\) is the ordered product Čech coordinate used
throughout this note; no further sign is suppressed.

The same resolution on \(K\) has zero terms in degrees0 and1 and
\(\chi_{\rm dil}k\) in degree2. Therefore \[
\boxed{R\Gamma(Y^2,K)=\chi_{\rm dil}k[-2],\qquad
H^2(Y^2,K)=\chi_{\rm dil}k.}
\tag{RPD3.3}
\] Its trace sends the stated top Čech coordinate to itself. The induced
cochain map of RPD2.3 is exactly \[
\beta_C^0=0,\quad\beta_C^1=0,\quad
\beta_C^2(a\otimes b)=B_A(a,b).
\tag{RPD3.4}
\] It is a cochain map because the two terms in RPD3.2 belong to
\(A\otimes J\) and \(J\otimes A\). This proof gives the degree2 and
every sign from the original ordered complex.

The entire algebraic cohomology of \(C\) is retained: \[
H^0(C)=H\otimes H,\quad
H^1(C)=(H\otimes Q)\oplus(Q\otimes H),\quad
H^2(C)=Q\otimes Q,\quad H=\ker d.
\tag{RPD3.5}
\] For a direct proof, split the vector-space rows
\(0\to H\to P\to J\to0\) and \(0\to J\to A\to Q\to0\) algebraically.
This decomposes \(D\) into \(H[0]\), \(Q[-1]\), and the contractible
complex \([J\xrightarrow{\mathrm{id}}J]\). Tensoring its explicit
identity homotopy, with the Koszul sign, contracts every summand
containing this last complex. The resulting cohomology is RPD3.5; its
isomorphisms are the canonical products of cohomology classes, so the
statement is independent of the chosen splittings. A continuous section
\(Q\to A\) is not asserted.

In degree2 one can avoid any splitting: the image of \(d_C^1\) is
exactly \(J\otimes A+A\otimes J\), since the two variables in RPD3.2
range independently. Its quotient is exactly \(Q\otimes Q\). The
resulting degree-two map is \[
\boxed{\overline\beta:Q\otimes Q\to\chi_{\rm dil}k,\qquad
q_1\otimes q_2\mapsto B_\zeta(q_1,q_2).}
\tag{RPD3.6}
\] Every endpoint and faithful extra copy remains in \(H\) and hence in
the first two groups of RPD3.5. The trace has zero components in those
degrees; their groups have not been deleted.

\subsection{RPD4. Currying gives exactly the existing global residue
map}\label{rpd4.-currying-gives-exactly-the-existing-global-residue-map}

Let \(L=\chi_{\rm dil}k[-2]\). Its only nonzero term is
\(L^2=\chi_{\rm dil}k\). Currying RPD3.4 gives the continuous-Hom
complex map \[
\Psi:D\longrightarrow\operatorname{Hom}_c(D,L)
=\chi_{\rm dil}D'_c[-2].
\tag{RPD4.1}
\] The nonzero terms of this target are \(A'\) in degree1 and \(P'\) in
degree2. With the usual Hom differential
\(d_{\rm Hom}f=d_Lf-(-1)^{\deg f}f d_D\), its differential from degree1
is \(+d'\). The complete curried map is \[
\boxed{\Psi^0=0,\qquad
\Psi^1(a)(b)=B_A(a,b),\qquad
\Psi^1=\pi'D_\zeta\pi.}
\tag{RPD4.2}
\] There is no additional minus sign in this currying: evaluation is
\(f\otimes b\mapsto f(b)\), and the Koszul signs are already those of
the tensor and Hom differentials.

Direct verification keeps both equations: \(\Psi^1(dx)(b)=B_A(dx,b)=0\),
and \(d'\Psi^1(a)(x)=B_A(a,dx)=0\). These are precisely the two sheaf
naturality conditions in RPD2, now evaluated as the global cochain
equations. The induced degree-one cohomology map is the existing
injection \(D_\zeta:Q\to\chi_{\rm dil}Q'\).

Its continuity into the actual strong dual follows from RPD1.5. For
every bounded set \(T\subset A\), its dual seminorm is bounded by \[
\sup_{b\in T}|B_A(a,b)|
\le\frac{\zeta(2)(1+4\pi^2)}{\pi}
\,b_{2,1}(\mathcal M_0a)\,
\sup_{b\in T}b_{2,1}(\mathcal M_0b).
\] The last supremum is finite by continuity of \(\mathcal M_0\). Thus
this currying lands in the actual continuous dual, rather than requiring
an identification with the full algebraic dual.

Finally \[
\Psi(T_aa)(b)=B_A(T_aa,b)
=a\,B_A(a,T_{a^{-1}}b)
\] gives exactly the character-twisted contragredient action in ASD.
This proves the product-site bridge to its original global \(\Psi\),
without contradicting FTD9's different same-site vanishing statement.

\subsection{RPD5. The actual diagonal topology and sheaf
pullback}\label{rpd5.-the-actual-diagonal-topology-and-sheaf-pullback}

Let \(\Delta:Y\to Y^2\), \(x\mapsto(x,x)\). It is continuous because
\(\Delta^{-1}(U\times V)=U\cap V\); its inverse onto its image is either
coordinate projection, so it is a topological embedding.

It is not closed. The point \((\eta,\eta)\) is in the closure relation
above every point: every nonempty basic open contains it. Hence its
closure is all of \(Y^2\). The diagonal contains this point but omits,
for example, \((c_+,c_-)\), so its closure is larger than itself. It is
not open either: the smallest neighbourhood of \((c_+,c_+)\) contains
\((c_+,\eta)\), which is not diagonal. Because a dense locally closed
subset would be open, the diagonal is not locally closed.

Inverse-image sheaves on these finite diagrams have the actual stalks at
their image points. Consequently \[
\Delta^{-1}(F\boxtimes F)=F\otimes F
=(W_+\otimes W_+\xrightarrow{r_+\otimes r_+}A\otimes A
\xleftarrow{r_-\otimes r_-}W_-\otimes W_-),
\] \[
\Delta^{-1}K=j_{\eta!}\chi_{\rm dil}k.
\tag{RPD5.1}
\] The map \(\Delta^{-1}\beta\) is \(B_A\) at \(\eta\) and zero on the
two closed stalks. It remains a nonzero sheaf morphism. Its naturality
now requires vanishing on \(J\otimes J\), which follows from the
stronger two-slot vanishing RPD1.6.

The original ordered one-dimensional Čech complex gives \[
R\Gamma(Y,j_{\eta!}\chi_{\rm dil}k)
=\chi_{\rm dil}k[-1].
\tag{RPD5.2}
\] In particular its \(H^1\) is \(\chi_{\rm dil}k\), whereas its \(H^2\)
is zero. The pullback on sheaves is exact over \(k\), but it has not
been identified with a closed-diagonal Gysin map; the topology just
computed precludes that identification without another constructed
operation.

\subsection{RPD6. The precise global diagonal comparison and its
signs}\label{rpd6.-the-precise-global-diagonal-comparison-and-its-signs}

A concrete diagonal Čech comparison from \(C=D\otimes D\) to the Čech
complex \(C_\Delta\) of RPD5.1 has degree0 \[
\kappa^0(x\otimes y)=(x_+\otimes y_+,\,x_-\otimes y_-),
\] degree1 \[
\kappa^1(a\otimes y,\ x\otimes b)
=a\otimes r_-(y_-)+r_+(x_+)\otimes b,
\qquad \kappa^2=0.
\tag{RPD6.1}
\] Here \(x=(x_+,x_-)\), \(y=(y_+,y_-)\), and \[
d_\Delta(z_+,z_-)
=(r_+\otimes r_+)z_+-(r_-\otimes r_-)z_-.
\] The degree-zero cochain equation is the exact identity \[
\begin{aligned}
\kappa^1d_C(x\otimes y)
&=(r_+x_+-r_-x_-)\otimes r_-y_-
+r_+x_+\otimes(r_+y_+-r_-y_-)\\
&=r_+x_+\otimes r_+y_+-r_-x_-\otimes r_-y_-\\
&=d_\Delta\kappa^0(x\otimes y).
\end{aligned}
\tag{RPD6.2}
\] The remaining cochain equation has both sides zero. The comparison is
natural in the product sheaf: for a general product-site sheaf, replace
each displayed tensor restriction by its actual restriction along the
corresponding rectangle. In degree0 it is exactly restriction of
compatible global sections to the diagonal. The source complex comes
from the projective resolution in RPD3, so its cohomology is the
universal derived global-section functor. Therefore this natural cochain
comparison extending degree-zero restriction induces the actual global
pullback maps on every cohomology group.

Interchanging the choices of plus and minus in the degree-one formula
gives another such comparison. Their difference in degree1 is
\(-a\otimes dy+dx\otimes b=d_C^1(a\otimes y,x\otimes b)\). The map
\(h:C^2\to C_\Delta^1\) equal to identity, zero elsewhere, is the
explicit cochain homotopy between them. Thus the prescribed orientation
is transparent and the induced cohomology map is unchanged.

For the target \(K\), its global product complex is concentrated in
degree2 and its diagonal complex in degree1. The degree-preserving
diagonal comparison is the zero map. Moreover the source-target square
commutes strictly: composing RPD6.1 with \(B_A\) in degree1 gives \[
B_A(a,r_-y_-)+B_A(r_+x_+,b)=0
\tag{RPD6.3}
\] by the two original annihilation identities. All other possible
components are zero. Thus the nonzero degree2 product trace pulls back
to zero, while the nonzero diagonal sheaf pairing occurs in degree1.
These statements describe different degrees of the actual comparison,
rather than contradicting one another.

The lower-degree groups of the product are also accounted for. Under
RPD3.5 the map on \(H^0\) is the diagonal product of global sections. A
representative of \(H\otimes Q\) is \(x\otimes b\) with \(dx=0\), and
its diagonal image is \([r_+x_+\otimes b]\in H^1(C_\Delta)\), with
\(r_+x_+=r_-x_-\in J\). A representative of \(Q\otimes H\) is
\(a\otimes y\), and its image is \([a\otimes r_-y_-]\). Thus both images
lie in the mixed \(J\)-terms calculated next. Since the global Fourier
graph maps onto \(J\), these two image spaces together give every such
mixed term. Endpoints and extra copies have zero restriction and
therefore map to zero in this degree, while remaining in \(H\) itself.

\subsection{RPD7. The full diagonal quotient and both distinct
kernels}\label{rpd7.-the-full-diagonal-quotient-and-both-distinct-kernels}

Algebraic tensor over \(k\) preserves surjections. Therefore the two
diagonal restriction images are both \(J\otimes J\), and \[
\boxed{H^1(C_\Delta)=Q_\Delta:=(A\otimes A)/(J\otimes J).}
\tag{RPD7.1}
\] Its degree-zero cohomology is the complete fibre product of
\(W_+\otimes W_+\to J\otimes J\leftarrow W_-\otimes W_-\). If
\(E_\pm=\ker r_\pm=\mathbb C^2_\pm\oplus V_\pm^{\rm extra}\), the kernel
of each of these tensor restrictions is \[
K_\pm^\Delta=
E_\pm\otimes W_\pm+W_\pm\otimes E_\pm.
\] The intersection of its two indicated summands is
\(E_\pm\otimes E_\pm\), as follows by taking a vector-space complement
to \(E_\pm\). Consequently the complete degree-zero row is \[
0\to K_+^\Delta\oplus K_-^\Delta
\to H^0(C_\Delta)\to J\otimes J\to0.
\tag{RPD7.2}
\] The last map takes the common restriction and is onto because both
restriction maps are onto. This retains all diagonal endpoint and extra
contributions.

There is a canonical surjection \[
\vartheta:Q_\Delta\to Q\otimes Q,
\qquad [a\otimes b]\mapsto\pi a\otimes\pi b.
\tag{RPD7.3}
\] Its exact kernel is \[
\boxed{
0\to(J\otimes Q)\oplus(Q\otimes J)
\longrightarrow Q_\Delta\xrightarrow{\vartheta}Q\otimes Q\to0.}
\tag{RPD7.4}
\] For the first injection, \(j\otimes[b]\) maps to \([j\otimes b]\) and
\([a]\otimes j\) maps to \([a\otimes j]\). A different lift changes the
tensor by \(J\otimes J\), so these maps are well-defined. In
\(A\otimes A\), \((J\otimes A)\cap(A\otimes J)=J\otimes J\), verified by
a vector-space complement of \(J\). Their quotients modulo this
intersection are respectively \(J\otimes Q\) and \(Q\otimes J\). This
proves injectivity, the direct-sum assertion and exactness. It also
proves the image description at the end of RPD6.

The scalar diagonal pairing is \[
\beta_\Delta:Q_\Delta\xrightarrow{\vartheta}Q\otimes Q
\xrightarrow{\overline\beta}\chi_{\rm dil}k.
\tag{RPD7.5}
\] The kernel in RPD7.4 is the kernel of \(\vartheta\), not the whole
kernel of this scalar map. Its complete scalar kernel has the exact row
\[
\boxed{
0\to(J\otimes Q)\oplus(Q\otimes J)
\to\ker\beta_\Delta\to\ker\overline\beta\to0.}
\tag{RPD7.6}
\] This follows by restricting the surjection \(\vartheta\) to the
preimage of \(\ker\overline\beta\); every element of that scalar kernel
has a lift, and the kernel of the restriction remains RPD7.4.

For an entirely explicit scalar splitting, choose any actual nontrivial
zero \(\rho\), of order \(m\), and the entire Mellin jet isolators
\(E_{\rho,j}\) constructed in GZR. Define their unshifted quotient
representatives by \[
e_{\rho,j}=\pi\bigl(M_0^{-1}E_{\rho,j}\bigr)\in Q.
\] The inverse exists in the original test space by the exact Mellin
return. These are not \(\mathcal T k_{\rho,j}\): the latter has Mellin
transform \(2E_{\rho,j}\), so using it in both slots would multiply the
following value by four. GZR's full residue calculation with the
specified \(M_0^{-1}\) representatives gives \[
B_\zeta(e_{\rho,m-1},e_{1-\rho,0})
=\frac{m!}{\zeta^{(m)}(\rho)}.
\] Therefore \[
z_0=e_{\rho,m-1}\otimes
\frac{\zeta^{(m)}(\rho)}{m!}e_{1-\rho,0}
\quad\text{satisfies}\quad\overline\beta(z_0)=1.
\tag{RPD7.7}
\] The linear map \(z\mapsto z-\overline\beta(z)z_0\) projects onto the
entire kernel \(\ker\overline\beta\). This is a chosen scalar splitting,
not a claim of group-equivariant splitting. Nondegeneracy of \(B_\zeta\)
in its two slots does not make \(\overline\beta:Q\otimes Q\to k\)
injective. For example, isolated nonzero primary tensors whose second
character is not the reflection of the first have zero pairing by GZR6,
while remaining nonzero tensors. All such terms remain in RPD7.6.

\subsection{RPD8. The actual diagonal image of the product and the
scalar
defect}\label{rpd8.-the-actual-diagonal-image-of-the-product-and-the-scalar-defect}

Combining RPD6--RPD7 gives the exact mathematical relation between the
two global calculations: \[
H^1(Y^2,F\boxtimes F)
=(H\otimes Q)\oplus(Q\otimes H)
\longrightarrow Q_\Delta
\] has image precisely \((J\otimes Q)\oplus(Q\otimes J)=\ker\vartheta\).
Its cochain representatives were computed in RPD6. Thus the diagonal
scalar pairing kills the entire image of global product degree1. The
independent quotient \(Q\otimes Q\) on which it is nonzero appears in
product degree2, but that degree pulls back to zero because the diagonal
Čech dimension is1.

This establishes more than a mismatch of dimensions: the natural
comparison map, its image, the remaining quotient, and the full scalar
kernel have been explicitly constructed. It neither removes the product
pairing nor pretends that its trace can be transferred by a
closed-diagonal pullback.

\subsection{RPD9. Full multiplicity, companion signs and the exact
weight
content}\label{rpd9.-full-multiplicity-companion-signs-and-the-exact-weight-content}

At an actual zero \(\rho\) of order \(m\), retain \[
\zeta(\rho+t)=t^m u_\rho(t),\qquad
u_\rho(t)=\sum_{k\ge0}\frac{\zeta^{(m+k)}(\rho)}{(m+k)!}t^k,
\quad
\frac1{u_\rho(t)}=\sum_{k\ge0}c_{\rho,k}t^k.
\] The original local residue form is \[
\sum_{i+j+k=m-1}
\frac{F^{(i)}(\rho)}{i!}
\frac{G^{(j)}(1-\rho)}{j!}
(-1)^j c_{\rho,k}.
\tag{RPD9.1}
\] All contributing original derivatives and nilpotents remain in the
product morphism. GZR's global primary adjoint relation says that it
pairs \(Q_\rho\) only with \(Q_{1-\rho}\). Their exact original actions
are \[
T_a|_{Q_\rho}=a^\rho
\sum_{j=0}^{m-1}\frac{(\log a)^j}{j!}N_\rho^j,\qquad
T_a|_{Q_{1-\rho}}=a^{1-\rho}
\sum_{j=0}^{m-1}\frac{(\log a)^j}{j!}N_{1-\rho}^j.
\tag{RPD9.2}
\] The residue form satisfies \[
B_\zeta(N_\rho x,y)=-B_\zeta(x,N_{1-\rho}y).
\tag{RPD9.3}
\] One verifies this by multiplying \(F(\rho+t)\) by \(t\) and observing
that the second reflected argument introduces \(-t\). Hence the two
finite exponential nilpotents cancel in the simultaneous action and the
full character is \(a^{\rho+1-\rho}=a\). This proves the covariance
without discarding any Jordan level.

At every prime \(p\), the numerical character weights of the two factors
are \(2\Re\rho\) and \(2(1-\Re\rho)\), whose sum is2. The trace target
\(\chi_{\rm dil}\) has weight2. The proved equation fixes the sum, not
either factor separately. All formulas hold for the actual zero
parameter without setting its real part to \(1/2\), so this product
realization alone gives no additional weight separation. This is a
precise statement about what its exact equations contain, not a claim
that an off-critical zero exists.

For the companion denominator \(\zeta(1-s)\), GZR proves \[
B_\zeta(y,x)=-B_{\zeta^\vee}(x,y).
\tag{RPD9.4}
\] The factor swap on the tensor complex is
\(x\otimes y\mapsto(-1)^{\deg x\deg y}y\otimes x\); on its degree2 part
this is negative interchange. The geometric coordinate interchange also
acts by minus identity on the ordered degree2 trace, because the two
Čech directions are interchanged. These signs remain distinct from the
companion sign in RPD9.4. No symmetry, Hermitian form or positivity is
inferred for \(B_\zeta\) itself.

\subsection{RPD10. Source actions, labels and the completed forward
step}\label{rpd10.-source-actions-labels-and-the-completed-forward-step}

On the product use the tensor of the receiving rings. Its action on the
generic tensor is by the product of the two arithmetic scalar
coordinates; the target scalar representation uses that same product. If
either argument is in a source-faithful extra closed copy, its
restriction to the generic point is zero, as prescribed by the original
sheaf. The product sheaf retains every such copy. This does not identify
source \(\tau\) with integer1: their actions still differ on those
copies. In particular simultaneous source integer action on two inputs
multiplies the pairing by \(n^2\), whereas simultaneous original
spectral dilation at \(a=n>0\) gives the separate factor \(n\) of
RPD1.9. These are different specified operations.

For the established support carriers, a bilinear amplitude map retains
the ordered pair of input labels: \[
((x,\lambda),(y,\mu))\longmapsto
(B_\zeta(x,y),(\lambda,\mu)).
\] A nonzero amplitude requires both top labels; a zero keeps the full
pair. Linear comparison maps preserve each existing label by
\((v,\lambda)\mapsto(f(v),\lambda)\). Independent chart and tensor
factors keep their separate labels; no supported zero becomes primitive
\(\tau\).

The forward calculation after FTD's same-site nonlift is therefore
complete: the product-site sheaf morphism exists, its ordered degree2
trace gives exactly the original global residue pairing, and currying
gives the actual continuous-dual \(\Psi\). Its diagonal pullback is
explicitly computed on sheaves and on every Čech degree. The diagonal is
not closed, and the natural global pullback kills the degree2 trace
while retaining a distinct degree1 scalar map with kernel RPD7.6. The
full multiplicity equations fix the sum of paired numerical weights;
they do not establish purity of either factor. This supplies the precise
product and diagonal objects for the next geometric calculation, without
replacing that calculation by an assumed weight bound.

\subsubsection{The derived diagonal operation restores the original
trace}\label{the-derived-diagonal-operation-restores-the-original-trace}

The next calculation is completed in DIAGONAL\_EXTRAORDINARY\_RETURN.md.
Define q:Y\^{}2→Y by q(c+,c+)=c+, q(c-,c-)=c-, and send the other seven
points to eta. Its inverse images of the three minimal opens are
computed explicitly. The stalk identity Delta\_\emph{=q\^{}\{-1\},
including all restrictions, proves that the derived right adjoint of
Delta\_} is Rq\_*. It is not ordinary diagonal pullback.

For K=j\_(eta,eta)!chi C, the full injective complex has terms I\_eta,
I\_eta\^{}4, and I\_+ direct-sum I\_eta\^{}2 direct-sum I\_- in
degrees0,1,2. All signs and contractions are retained in the proof. Its
cohomology is j\_eta!chi C{[}-1{]}, and derived global sections recover
the original chi C{[}-2{]} trace with coefficient +1. Thus the degree
loss under ordinary pullback is not a general obstruction to a
diagonal-derived realization. This follows by the constructed right
adjoint, not by an assumed closed immersion or a numerical purity
statement. The full global identity RΓ Rq\_*=RΓ also retains every other
original cohomological degree.

\subsubsection{Full arithmetic correspondence and its derived
coefficient
return}\label{full-arithmetic-correspondence-and-its-derived-coefficient-return}

DERIVED\_RETURN\_FULL\_SOURCE\_CORRESPONDENCE.md, FSC0--FSC6, returns
the receiving construction to the entire original arithmetic source.
FSC2 rules out precisely a continuous single-valued lift of q that fixes
the complete arithmetic diagonal. FSC3 immediately constructs the full
alternative Z=\{(x,y,z):fz=q(fx,fy)\}, with projections a,b and lifted
diagonal x↦(x,x,x). Each nonexceptional a-fibre is the whole Spec Z; the
two equal closed-corner fibres are singletons.

FSC6 proves the actual canonical derived comparison
Ra\_\emph{b\textsuperscript{\{-1\}f}\{-1\}G ≅ (f×f)\^{}\{-1\}Delta\_}G
for bounded complexes of algebraic sheaves of complex vector spaces on
Y. It computes the inverse images of all three receiving injectives,
their exact direct images, and a functorial injective resolution; no
general proper-base-change theorem is assumed. The comparison applies to
the full derived input and trace target from DER and to their morphism.
Every arithmetic point remains in the correspondence. The coefficient
comparison does not assign a numerical weight to primitive Z\_1/tau or
settle the active numerical weight-separation target.

\subsubsection{The degree-one mixed obstruction has an explicit global
equivariant
contraction}\label{the-degree-one-mixed-obstruction-has-an-explicit-global-equivariant-contraction}

GLOBAL\_MIXED\_RETURN\_CONTRACTION.md, GMC0--GMC7, follows the generic
mixed row into its complete global sheaf complex. The actual section
s\_+(j)=(Sigma\^{}\{-1\}j,0,0) is continuous and dilation equivariant,
using the already proved summation inverse. It gives
S(m\_1,m\_2)=(-(s\_+ tensor1)m\_1,(1 tensor s\_+)m\_2,0,0), with d\_M
S=id. The exact global cochain map is (id-Sd\_M,vartheta), onto K\_0 in
degree0 and Q tensor Q in degree1, with zero differential. Its entire
kernel is {[}S(M)→iota(M){]}, contracted by h(iota m)=Sm. Every original
prime action intertwines.

Thus the degree-one mixed obstruction is killed by the constructed
contraction of ker(F)={[}S(M)→iota(M){]}. The full mixed subsheaf
retains RΓ(Y,N)≃K\_0{[}0{]}, with K\_0=(H tensor Q) direct-sum (Q tensor
H); it is not asserted to be acyclic. The generic polynomial
representative P(L\_total)z=iota m becomes the exact global boundary
d(Sm). All original endpoint and extra copies remain in K\_0=(H tensor
Q) direct-sum (Q tensor H); the original residue trace factors through
the computed map with coefficient +1. This is a theorem in the specified
algebraic tensor/sheaf calculation, with no assumption of
completed-tensor cohomology, mirror equivariance of the plus-only
section, or numerical purity. The remaining Q operators and all primary
nilpotents are retained explicitly in GMC6.3.

\subsection{Publication source and dependency
links}\label{publication-source-and-dependency-links}

The source geometry is Alain Connes and Caterina Consani, \emph{Schemes
over F1 and zeta functions},
\href{https://arxiv.org/abs/0903.2024v3}{arXiv:0903.2024v3, §5}. The
weight-control comparison is Pierre Deligne, \emph{La conjecture de
Weil. II}, Publications mathematiques de l'IHES52 (1980),137-252,
\href{https://www.numdam.org/item/PMIHES_1980__52__137_0/}{§§3.3.11
and3.6}. Deligne material was read through the identified French
transcription and separately recorded peer source-page checks, not
Deligne-authored TeX. These sources are not asserted to contain the new
programme derivations. The
\href{https://github.com/KokunoYumeto/zeta-function-research-reader/blob/36a82ecca98addb8a98b48204e349737856c7223/workbenches/splitzero-tandem/continuations/20260924-cohomology-weight-and-character-lifts/BUILD_AND_REVIEW.md}{preceding
complete proof and reading record} records inherited source versions and
actual inspection limits. The investigator's corrected construction
remains attributed in the original text above.
